\writeSectionFrame{\presentationTitle}{Das Problem}
\begin{frame}{Bilder in einem Dokument}
	\begin{center}
 		\includegraphics[scale=0.35]{\BILDERPFAD/problem1.png}
 	\end{center}
\end{frame}

\note{
 	Wir stellen uns einen Dokumenteditor vor, der die Möglichkeit bietet, Bilder in einem
 	Dokument zu speichern. Bilder können sehr groß sein und Dokumente ebenfalls und so kann
 	es sehr lange dauern, ein großes Dokument mit vielen Bildern zu öffnen. Das Öffnen eines
 	Dokuments sollte jedoch nicht lange dauern und deswegen sollten nicht alle Bilder sofort
 	geöffnet werden. Vor allem, weil man niemals das gesamte Dokument auf einmal sieht. Wie
 	können wir das erreichen, ohne den Code komplizierter zu machen?
}


\begin{frame}{Proxy vor Bildobjekten}
	\begin{center}
 		\includegraphics[scale=0.3]{\BILDERPFAD/problem2.png}
 	\end{center}
\end{frame}

\note{
 	Als Lösung betten wir nicht direkt die Bilder in das Dokument ein, sondern nur Proxy-Objekte. Der
 	Proxy ist ein Platzhalter und das Bild wird bei Bedarf geladen. Der Proxy verhält sich genauso wie
 	das Bild und sorgt dafür, dass das Bild erzeugt wird, wenn es benötigt wird. 
}

\begin{frame}{Beispiel aus der realen Welt}
	\begin{center}
 		\includegraphics[scale=0.25]{\BILDERPFAD/problem3.png}
 	\end{center}
\end{frame}

\note{
 	Aus der realen Welt kennen wir das Prinzip, wenn wir mit Karte bezahlen. Die Karte ist im Geschäft
 	ein Proxy für Bargeld. Wir benutzen es prinzipiell wie Bargeld, nur müssen wir nicht Bargeld mit
 	uns herumschleppen. Der Betrag wird dann auf das Konto des Händlers gebucht und auch dessen Konto
 	ist ein Proxy.  
}